\documentclass[10pt,a4paper]{article}
\usepackage[margin=2.2cm]{geometry}
\usepackage{amsmath,amssymb}
\usepackage{microtype}
\usepackage[colorlinks=true,linkcolor=blue,urlcolor=blue,citecolor=blue]{hyperref}

\newcommand{\Hk}{\mathcal H_k}
\newcommand{\R}{\mathbb R}
\newcommand{\E}{\mathbb E}

\title{From RKHS Interpolation to Learning:\\Kernel Ridge Regression and Gaussian Processes}
\author{Eldad Kronfeld}
\date{}
% Draft: assignment requirements and lecture-specific alignment remain unverified.
\begin{document}
\maketitle
\vspace{-2em}

\begin{abstract}
	RKHS Interpolation connect the pointwise measurements to the inner product structure using reproducing kernels, in here I will show how to deal with a "relaxed" optimization problem in a regime with noisy measurements.
	I will do that by establishing a new \textbf{Optimization Problem} (OP), which will be solved with a technique knowen as \textbf{Kernel Ridge Regression} (KRR)
	where we use the geometry of the RKHS to produce a minimizer.
\end{abstract}

\textbf{Remark:} In case you want to checkout the background work, it is located in: \href{https://github.com/Eldadkro/manifold-approx-solutions}{repository}.
\section{RKHS interpolation: exact agreement with the data}
Let $\mathcal X\subseteq\R^d$ be an input domain and let $k:\mathcal X\times\mathcal X\to\R$ be a symmetric positive-semidefinite kernel: for every finite collection $z_1,\ldots,z_m\in\mathcal X$ and $a\in\R^m$,
$\sum_{i,j=1}^m a_i a_j k(z_i,z_j)\geq0$.
Its associated RKHS $\Hk$ is a Hilbert space of functions on $\mathcal X$ with
\begin{equation}
	k(\cdot,x)\in\Hk,\qquad f(x)=\langle f,k(\cdot,x)\rangle_{\Hk}.
	\label{eq:reproduction}
\end{equation}
In particular, evaluation is bounded: $|f(x)|\leq\|f\|_{\Hk}\sqrt{k(x,x)}$.

Fix inputs $x_1,\ldots,x_n\in\mathcal X$ and observed responses $y=(y_1,\ldots,y_n)^T\in\R^n$. Define the sample-dependent Gram matrix and kernel vector by
\[
	K_{ij}=k(x_i,x_j),\quad K\in\R^{n\times n},\qquad
	k_x=(k(x_1,x),\ldots,k(x_n,x))^T\in\R^n.
\]
The kernel is a function on pairs of inputs; $K$ is its evaluation on this particular sample. \\
For now assume $K$ is invertible, so every response vector can be \textbf{interpolated}.
\\
\textbf{Remak:} After already researching KRR, I somewhat regret that I didn't think to research the problems in the OP of the Optimal intarpolant...
would have been a great single page work,
but alas I already invested too much so I will
only assume that the Gram matrix isn't singular and that the feasible set isn't empty.
\subsection{Optimal interpolant in reproducing kernel Hilbert space}
Given $\Hk$ RKHS we wish to find the optimal solution to the following OP, such that the esimator (interpolant) fits the measurements exactly and is the one with the smallest possible norm.
\begin{equation}
	f_{\mathrm{int}}=\operatorname*{arg\,min}_{f\in\Hk}\|f\|_{\Hk}^2
	\quad\text{subject to}\quad f(x_i)=y_i\quad(i=1,\ldots,n).
	\label{eq:interpolation}
\end{equation}

The reproducing property reduces this problem to finite dimensions. Set
$V=\operatorname{span}\{k(\cdot,x_1),\ldots,k(\cdot,x_n)\}$ and decompose $f=f_V+f_\perp$ orthogonally. Equation~\eqref{eq:reproduction} gives $f_\perp(x_i)=0$, whereas
\begin{equation}
	\|f\|_{\Hk}^2=\|f_V\|_{\Hk}^2+\|f_\perp\|_{\Hk}^2.
	\label{eq:split}
\end{equation}
Removing $f_\perp$ preserves the observed values and decreases the norm unless $f_\perp=0$. Thus the minimum-norm interpolant lies in $V$. Writing it as $\sum_{j=1}^n\alpha_j k(\cdot,x_j)$, the constraints become $K\alpha=y$, yielding
\begin{equation}
	\boxed{f_{\mathrm{int}}(x)=k_x^TK^{-1}y.}
	\label{eq:interpolant}
\end{equation}
Interpolation fixes the training values; the RKHS norm determines how they are extended to unobserved inputs. This preference depends on the kernel, rather than on a universal notion of smoothness.
Because of Moore–Aronszajn.

\section{Kernel ridge regression: relaxing interpolation}
Exact agreement with the measured values may be undesirable when responses contain measurement noise. Instead of enforcing the constraints in~\eqref{eq:interpolation}, KRR penalizes their violation:
\begin{equation}
	\widehat f_\lambda=\operatorname*{arg\,min}_{f\in\Hk}
	\left\{\frac1n\sum_{i=1}^n(y_i-f(x_i))^2+\lambda\|f\|_{\Hk}^2\right\},
	\qquad\lambda>0.
	\label{eq:krr}
\end{equation}
The first term measures training error; the second controls the RKHS norm, which in this regime could be interpreted as the preference for a regressors of a
spacific form aside
from the inductive bias we introduced by picking the kernel and "locking" the solutions into the spacific RKHS.\\
This is a learning rule because it maps the observed sample to a predictor at new inputs. Regularization expresses a preference that can help with noisy data,
but does not by itself guarantee generalization.

The same orthogonal decomposition applies: removing $f_\perp$ leaves the loss unchanged because the ortogonal
element doesn't effect the measured points (first part of the loss) becuase:
\[
	\text{for: } h\in f_\perp,i=1,2..,n \qquad \langle h,k(\cdot,x_i)\rangle_{\Hk} = 0.
\]
and strictly punish the target function if $f_\perp\ne0$. Hence $\widehat f_\lambda\in V$. This is the representer argument for this particular objective. For $f=\sum_j\alpha_jk(\cdot,x_j)$,
\[
	(f(x_i))_{i=1}^n=K\alpha,\qquad \|f\|_{\Hk}^2=\alpha^TK\alpha,
\]
so the coefficient objective is
$J(\alpha)=n^{-1}\|y-K\alpha\|_2^2+\lambda\alpha^TK\alpha$.
Its stationarity equation is
\begin{equation}
	\nabla J(\alpha)=\frac2n K\bigl[(K+n\lambda I_n)\alpha-y\bigr]=0.
\end{equation}
We need not cancel $K$, which might be singular. Since $K$ is positive semidefinite and $\lambda>0$, $K+n\lambda I_n$ is positive definite. The vector
$\alpha_\lambda=(K+n\lambda I_n)^{-1}y$ satisfies stationarity and minimizes the convex objective. Therefore
\begin{equation}
	\boxed{\widehat f_\lambda(x)=k_x^T(K+n\lambda I_n)^{-1}y.}
	\label{eq:krr-predictor}
\end{equation}
The function minimizer is unique because~\eqref{eq:krr} is strictly convex. If $K$ is singular, coefficient vectors need not be unique, but their differences in $\ker K$ represent the zero function. In computation, the inverse notation means solving a linear system.

The fitted training values satisfy
\[
	y-K\alpha_\lambda=n\lambda(K+n\lambda I_n)^{-1}y,
\]
so exact agreement is generally lost. When $K$ is invertible, letting $\lambda\downarrow0$ recovers~\eqref{eq:interpolant}, including convergence in $\Hk$. Thus KRR is a
relaxation of minimum-norm interpolation, not a separate construction. Unlike exact interpolation for arbitrary $y$, it remains well-defined for singular $K$.
\end{document}
